\chapter{Possibility Space Mapping}
\label{chap:Possibility_Space_Mapping}

\section{Purpose}

Possibility Space Mapping (PSM) is how SEGO:
\begin{itemize}
    \item organizes candidate trajectories,
    \item charts ``regions'' of behaviour or narrative,
    \item and makes the space of options navigable for EGO and HITL.
\end{itemize}

Rather than generating one option at a time, PSM produces a map of
what \emph{could} happen, under different assumptions.

\section{Space Representation}

We represent a possibility space $\mathcal{P}$ as:
\[
\mathcal{P} = \langle \mathcal{S}, \mathcal{T}, \mathcal{E} \rangle
\]
where:
\begin{itemize}
    \item $\mathcal{S}$: states (nodes),
    \item $\mathcal{T}$: transitions (edges),
    \item $\mathcal{E}$: evaluation metadata (scores, risks, tags).
\end{itemize}

Each state $s \in \mathcal{S}$ has:
\begin{itemize}
    \item a W5H annotation,
    \item Engram references,
    \item tag-based semantics.
\end{itemize}

\section{Mapping Process}

\begin{enumerate}
    \item Start from current context $q$ and resonant set $R(q)$.
    \item Use Manifold Expansion to generate candidate states and
          transitions.
    \item Label states with:
          \begin{itemize}
              \item W5H snapshots,
              \item key tags (e.g.\ ONTO, EPI, PHIL),
              \item alignment and risk markers.
          \end{itemize}
    \item Build a graph or lattice $\mathcal{P}$ that EGO can traverse.
\end{enumerate}

\section{Use by EGO and HITL}

\begin{itemize}
    \item EGO can:
          \begin{itemize}
              \item compare multiple trajectories at once,
              \item search for Pareto-optimal or dominated options,
              \item reason about tradeoffs explicitly.
          \end{itemize}
    \item HITL can:
          \begin{itemize}
              \item inspect high-risk branches,
              \item approve or forbid regions of the space,
              \item adjust Governance Policies based on observed structures.
          \end{itemize}
\end{itemize}

\section{Constraints}

\begin{itemize}
    \item PSM must respect drift and policy constraints when constructing
          and presenting options.
    \item It must clearly separate:
          \begin{itemize}
              \item explored-but-rejected paths,
              \item unexplored but plausible paths,
              \item committed history.
          \end{itemize}
\end{itemize}

\section{Summary}

Possibility Space Mapping turns raw generative potential into a legible
map of choices. It lets EGO and humans see not just \emph{what} options
exist, but how they relate and what they may cost.
